\section{Introduction}

Functions are one of the core abstractions that programmers use.
These functions - and their authors - have to communicate with different calling sites -
and their authors - in different contexts.
This communication is performed through the use of arguments.
However, the term "argument" is used in different ways depending on the speaker, listener,
and context.
For example, people who are working on a new function written in the C programming
language might use the word "arguments" to refer to the set of names that appear between
the parantheses in the function signature.
Those same people who are later debugging code which includes a call to that function
might use the word "arguments" to refer to the values which are given to the function.
This informal use of the word is sufficient for day-to-day use, but it starts to break
down when examined more closely.

Consider that the Rust programming lanugage allows programmers to specify arguments using
one of two mechanisms.
In the common case, programmers can use a plain name and type:

\begin{pygmented}[lang=rust]
fn jump_plain(starting: &Point) -> Point {
  Point {
    x: starting.x,
    y: starting.y * 2
  }
}
\end{pygmented}

It would not be controvertial to claim that this function contains one argument.
This argument happens to be associated with a local variable named \texttt{starting} which
is guaranteed to contain data of type \texttt{Point}.

When it is useful, programmers can also specify an argument by a pattern
\footnote{Technically the "plain name" is also a pattern, albeit a very simple one.}
and type:

\begin{pygmented}[lang=rust]
fn jump_pattern(Point {
                  x: horizontal,
                  y: vertical
                }: &Point)
  -> Point {

  Point {
    x: *horizontal,
    y: *vertical * 2
  }
}
\end{pygmented}

Does this version contain one argument or two?
We could say that it contains 2 arguments, \texttt{horizontal} and \texttt{vertical}.
In the context of the function definition this would make sense.
However, at the calling site both appear to accept only a single argument:

\begin{pygmented}[lang=rust]
jump_plain  (&uut);
jump_pattern(&uut);
\end{pygmented}

Consider also the conventions in shell programming languages.
Shell functions might include "options" which are arguments that may or may not be related
to the argument that follows it.
For example, if a function encounters an argument with the value "-{}-ignore" then it
might interpret the following argument as a specification of something to ignore.
But if it instead encounters an argument with the value "-{}-ignore=foo" it might
interpret "foo" as the thing to ignore.
The following argument would then be interpreted in an entirely different manner.
In common speech, it is easy to treat both "-{}-ignore" and "foo" as part of the same
argument: "the ignore argument has value foo".
When the both the name and value are provided in the same string ("-{}-ignore=foo"), this
is technically accurate from the perspective of the interpreter (it is a single argument).
However, when they are given in diffrent strings ("-{}-ignore" "foo") it is not (they are
two arguments which are semantically related).

The premise of this paper is to resolve the above inconsistencies by treating an
"argument" as a set of associations between concrete things.
In the Rust example both versions of the function have a single argument.
In both cases this argument is associated with a single value at the calling site.
However, in the pattern case the argument is associated with 2 local variables in the
function definition.
Neither the values at the calling sites nor the local variables in the definitions are the
arguments: the association between the values and the variables are the arguments.
In the Bash example, the "-{}-ignore" and "foo" values are just that: values.
When the function author uses a helper program like \texttt{getopt} to process these
values with a given argument specification, both of the values become associated with the
same argument.

This paper will explore the implications of the premise by describing the common and
unique features associated with function arguments across several languages at different
levels of abstraction.
Additionally, it will present a mechanism implemented in GNU Guile which frees
programmers to create arbitrary associations when specifying arguments.
This will generalize the features provided by various languages, allowing programmers to
use the set of features that make sense given the context they are working in.

\subsection{Reproducibility and Transparency}

The source code used in example snippets can be found in the online git repository located
at http://git.sr.ht/~skyvine/comparative-function-arguments.
The source code in this repository contains far more than the snippets provided in this
paper, because I checked a number of variations on things that turned out to be
uninteresting for the purpose of this paper, but I do not feel comfortable deleting the
information gathered from these attempts (the fact that the outcome is uninteresting is
itself interesting in some contexts).
The repository includes files appropriate for use with GNU Guix to reproduce the software
environment I used while creating this paper.
It pins to a specific revision of the main Guix channel so that updates do not interfere
with reproducibility.
The Makefile within the "examples" directory launches a pure shell for the user so that
environmental factors do not interfere with reproducibility.

The git repository also contains the bibtex file used to generate references in this
paper.
This source file contains additional information which is not present in the output, such
as the SHA256 of referenced tarballs and the commit hash of source repositories.
\footnote {
	I intend to incorporate this information into the references in the final version of
	this paper, but I have not used Latex before so I will need to learn how to do that.
}

\subsection{Clarifying Terminology}

As mentioned above, this paper defines the term "argument" to refer to a set of
associations between concrete code objects, such as values and variables.
Some languages provide similar features under different names and sometimes those names
do not align with the vocabulary used by this paper.
For example, different communities use different terms to refer to the values that shell
functions interpret as implying a value.
A shell function might recognize a value "-{}-verbose" to mean that a variable named
"verbose" should have the value "true".
I have heard this mechanism variously referred to as a flag, a switch, or an option.
This paper refers to this mechanism as an "implicit value" because the caller uses a name
associated with the argument to imply a value which does not appear explicitly.
A complete least of terms which describe a feature provided by more than one language
follows:

\textbf{Positional Value} \textit{(common name: N/A)}:
These values are associated with arguments based on the index of the value in the list of
all positional values.

\textbf{Default Value} \textit{(common name: Optional Argument)}:
A value that an author associates with an argument for use when
the caller declines to supply a value.

\textbf{Named Value} \textit{(common name: Keyword Argument, Option Argument
\footnote{
	The term "option argument" is not to be confused with the term "option\textbf{al}
	argument".
	The former is used by shell users to refer to an argument that come after an option
	(such as "-{}-ignore" "foo" in the previous example) while an option\textbf{al} argument
	is a term used by scripters and system programmers to refer to an argument associated
	with a default value.
})}:
These values are associated with an argument based on a name which the caller attaches to
the value.
The argument may be associated with several \textbf{synonymous} names.

\textbf{Implicit Value} \textit{(common names: Switch, Flag, Option)}:
An argument has an implicit value if the caller can specify a name associated with the
argument but omit the value.
For example, the ubiquitous "--help" flag in most CLI tools is an implicit value.
An argument with an implicit value may also have \textbf{antonyms}, which invert the
semantics of the primary name.

\textbf{Typed Value} \textit{(common name: N/A)}:
Some languages allow or require associating a type with an argument, in which case the
value provided by the caller must be compatible with that type.

\textbf{Destructuring} \textit{(common names: Pattern, Unpacking)}:
Some languages provide a mechanism to break a composite structure (which might be a struct
or a container such as a list) into its component parts.
In some cases this functionality is available to an author, who can declare that an
argument should have a specific type but the values within that type should be bound to
separate local variables.
In other cases it might be available to a caller, in which case the values within a
composite structure are associated with different arguments.
It is possible for a language to provide both of these variants simultaneously.

\textbf{Variadic Function} \textit{(common name: N/A)}:
A function which can accept an arbitrary number of values.
Variadic functions are distinct from arguments with default values because with default
values the author controls the range of acceptable argument counts.
With variadic functions, the caller can pass in an arbitrarily large number of arguments
(physical limitations notwithstanding).
Variadic values are values which are associated with this feature.
For example, a variadic function may require one positional value at the beginning of the
value list.
The first value is not a variadic value but all other values are.
A variadic variable is provided by some languages. This is a local variable bound to some
(possibly empty) aggregate datatype (typically a list) which contains variadic values.
